*   Role: Professional Mathematics teacher.
    *   Goal: Solve a specific LaTeX exercise (Thema Γ).
    *   Formatting: Strict adherence to the provided `FORMATTING_GUIDE.md`.
    *   Constraints:
        *   Environment: `\begin{thema}{Γ} ... \end{thema}`.
        *   Queries: `\begin{erwthma} ... \end{erwthma}`.
        *   Items: Use `\item`.
        *   Math: `$ ... $` for inline, `\[ ... \]` for display, `\begin{align*}` for multi-line.
        *   Terminology: `\eng{...}` for English terms (e.g., `\eng{Lagrange}`), `\textbf{...}` for emphasis.
        *   Integrals: Only definite integrals, use `\d x`.
        *   No indefinite integrals, no 2nd derivative test, no differential equations, no MVT for integrals, no differentiation of variable limit integrals.
        *   Specific symbols: `\hm`, `\syn`, `\ef`, `\syf` (though not needed here since there's no trig).
        *   No preamble, no instructions, no repetition of the problem statement.

    *   Given: $x > 1$, $(x-1)f'(x) - f(x) = (x-1)^2 e^x$, $f(2) = e^2$.
    *   Part 1: Divide by $(x-1)^2$ to show $f(x) = (x-1)e^x$.
        *   $(x-1)f'(x) - f(x) = (x-1)^2 e^x$
        *   Divide by $(x-1)^2$: $\frac{(x-1)f'(x) - f(x)}{(x-1)^2} = e^x$
        *   Recognize the LHS as the derivative of $\frac{f(x)}{x-1}$.
        *   Let $g(x) = \frac{f(x)}{x-1}$. Then $g'(x) = \frac{f'(x)(x-1) - f(x)(1)}{(x-1)^2}$.
        *   So $g'(x) = e^x$.
        *   Integrate from $2$ to $x$: $\int_2^x g'(t) \d t = \int_2^x e^t \d t$.
        *   $g(x) - g(2) = [e^t]_2^x = e^x - e^2$.
        *   $g(2) = \frac{f(2)}{2-1} = \frac{e^2}{1} = e^2$.
        *   $g(x) - e^2 = e^x - e^2 \Rightarrow g(x) = e^x$.
        *   $\frac{f(x)}{x-1} = e^x \Rightarrow f(x) = (x-1)e^x$.

    *   Part 2: Monotonicity, minimum at $x=1$, and range for $x \in [1, \infty)$.
        *   $f(x) = (x-1)e^x$.
        *   $f'(x) = 1 \cdot e^x + (x-1)e^x = e^x(1 + x - 1) = xe^x$.
        *   For $x > 1$, $x e^x > 0$.
        *   So $f$ is \textbf{γνησίως αύξουσα} on $(1, \infty)$.
        *   If we extend to $[1, \infty)$, $f(1) = (1-1)e^1 = 0$.
        *   Since $f$ is strictly increasing for $x > 1$ and continuous at $x=1$, the minimum value $y=0$ is attained at $x=1$.
        *   Range: $f(1)=0$, $\lim_{x \to \infty} f(x) = \lim_{x \to \infty} (x-1)e^x = \infty$.
        *   Range: $[0, \infty)$.

    *   Part 3: Use MVT (\eng{Lagrange}) on $[x, x+1]$ for $x > 1$ to prove $x e^{x+1} - (x-1)e^x > x e^x$.
        *   Apply \eng{Lagrange} to $f(x) = (x-1)e^x$ on $[x, x+1]$.
        *   $f$ is \textbf{συνεχώς} on $[x, x+1]$ and \textbf{παραγωγήσιμη} on $(x, x+1)$.
        *   $\exists \xi \in (x, x+1)$ such that $\frac{f(x+1) - f(x)}{(x+1) - x} = f'(\xi)$.
        *   $f(x+1) = (x+1-1)e^{x+1} = x e^{x+1}$.
        *   $f(x) = (x-1)e^x$.
        *   $f'(\xi) = \xi e^\xi$.
        *   So, $x e^{x+1} - (x-1)e^x = \xi e^\xi$.
        *   Since $x < \xi < x+1$ and $g(t) = t e^t$ is \textbf{γνησίως αύξουσα} for $t > 0$ (because $g'(t) = e^t(1+t) > 0$), we have $\xi e^\xi > x e^x$.
        *   Therefore, $x e^{x+1} - (x-1)e^x > x e^x$.

    *   Part 4: Area between $C_f$, tangent at $x=2$, and $x=1$.
        *   Point of tangency: $x=2, y=f(2)=e^2$.
        *   Slope of tangent: $f'(2) = 2e^2$.
        *   Equation of tangent $\epsilon$: $y - e^2 = 2e^2(x-2) \Rightarrow y = 2e^2 x - 4e^2 + e^2 \Rightarrow y = 2e^2 x - 3e^2 = e^2(2x-3)$.
        *   Area $E = \int_1^2 [f(x) - \epsilon(x)] \d x$ OR $\int_1^2 [\epsilon(x) - f(x)] \d x$.
        *   Check which is above: $f(x) = (x-1)e^x$.
        *   $f$ is convex for $x>1$? $f''(x) = e^x(1+x) > 0$ for $x>1$.
        *   Since $f$ is convex, the tangent lies below the curve.
        *   $E = \int_1^2 [(x-1)e^x - (2e^2 x - 3e^2)] \d x$.
        *   Calculate $\int_1^2 (x-1)e^x \d x$:
            *   Using parts: $\int u dv = uv - \int v du$. $u=x-1, dv=e^x dx \Rightarrow du=dx, v=e^x$.
            *   $\int (x-1)e^x \d x = (x-1)e^x - \int e^x \d x = (x-1)e^x - e^x = (x-2)e^x$.
            *   Definite integral: $[(x-2)e^x]_1^2 = (0)e^2 - (-1)e^1 = e$.
        *   Calculate $\int_1^2 (2e^2 x - 3e^2) \d x$:
            *   $\int_1^2 e^2(2x-3) \d x = e^2 [x^2 - 3x]_1^2 = e^2 [(4-6) - (1-3)] = e^2 [-2 - (-2)] = e^2 [0] = 0$.
            *   Alternatively: $\int_1^2 (2e^2 x - 3e^2) \d x$ is the integral of a linear function over $[1,2]$. The midpoint is $1.5$, and $2e^2(1.5) - 3e^2 = 0$. The area is 0.
        *   So $E = e - 0 = e$.

    *   Start with `%=========== ΛΥΣΗ - Xο ΘΕΜΑ Γ ===========`.
    *   `\begin{thema}{Γ}`.
    *   `\begin{erwthma}`.
    *   `\item` for each part.
    *   `\eng{Lagrange}`.
    *   `\textbf{...}` for monotonicity/continuity.
    *   `\[ ... \]` for display math.
    *   `\begin{align*}` for integration steps.
    *   `\d x` for differential.
    *   End with `%=========== ΤΕΛΟΣ ΛΥΣΗΣ - Χο ΘΕΜΑ Γ ===========`.

    *   *Part 1:* $\frac{(x-1)f'(x) - f(x)}{(x-1)^2} = e^x \implies \left(\frac{f(x)}{x-1}\right)' = e^x$. Correct.
    *   *Part 2:* $f'(x) = xe^x$. Correct. $f(1)=0$. Range $[0, \infty)$. Correct.
    *   *Part 3:* $f(x+1)-f(x) = f'(\xi)$. $\xi e^\xi > xe^x$. Correct.
    *   *Part 4:* $\epsilon(x) = 2e^2x - 3e^2$. $\int_1^2 f(x) \d x = e$. $\int_1^2 \epsilon(x) \d x = 0$. Result $e$. Correct.

    *Wait, check the integral $\int_1^2 (x-1)e^x \d x$:*
    $u=x-1, v'=e^x \Rightarrow u'=1, v=e^x$.
    $\int (x-1)e^x \d x = (x-1)e^x - \int e^x \d x = (x-1)e^x - e^x = (x-2)e^x$.
    Evaluating from 1 to 2: $(2-2)e^2 - (1-2)e^1 = 0 - (-e) = e$.
    Correct.

    *Check the tangent integral $\int_1^2 (2e^2 x - 3e^2) \d x$:*
    $e^2 [x^2 - 3x]_1^2 = e^2 ( (4-6) - (1-3) ) = e^2 (-2 + 2) = 0$.
    Correct.

    *Constraint Check:* No indefinite integrals? The prompt says "No indefinite integrals: use definite integrals with brackets $[...]_a^b$".
    So I should write:
    $\int_1^2 (x-1)e^x \d x = [(x-2)e^x]_1^2 = \dots$
    And $\int_1^2 e^2(2x-3) \d x = [e^2(x^2-3x)]_1^2 = \dots$

    *Terminology check:* "Squeeze Theorem" -> `\eng{Squeeze}`, "Mean Value Theorem" -> `\eng{Lagrange}`.

    *Final structure check:*
    `\begin{thema}{Γ}`
    `\begin{erwthma}`
    `\item` ...
    `\end{erwthma}`
    `\end{thema}`

    *Wait*, the prompt says "Provide ONLY the LaTeX code for the solution. Do NOT include a preamble."
    Also, start and end with specific markers.

    *Refining the integration in Part 4:*
    The function $f(x) = (x-1)e^x$ is convex on $[1, 2]$ because $f''(x) = (x+1)e^x > 0$.
    Therefore, the tangent line at $x=2$ lies below the curve on $[1, 2]$.
    The area is $\int_1^2 [f(x) - \epsilon(x)] \d x$.

    *Double check Part 3 inequality:*
    $x e^{x+1} - (x-1)e^x > x e^x$
    $f(x+1) - f(x) = f'(\xi) = \xi e^\xi$
    Since $\xi \in (x, x+1)$ and $x > 1$, then $\xi > x$.
    Since $g(t) = te^t$ is strictly increasing for $t > 0$, $\xi e^\xi > xe^x$.
    Thus $x e^{x+1} - (x-1)e^x > x e^x$.
    Correct.
%=========== ΛΥΣΗ - Xο ΘΕΜΑ Γ ===========
\begin{thema}{Γ}
\begin{erwthma}
\item Διαιρώντας τη δεδομένη σχέση $(x-1) f'(x) - f(x) = (x-1)^2 e^x$ με το $(x-1)^2$ για $x > 1$, προκύπτει:
\[ \frac{(x-1) f'(x) - f(x)}{(x-1)^2} = e^x \]
Παρατηρούμε ότι το αριστερό μέλος είναι η παράγωγος της συνάρτησης $g(x) = \frac{f(x)}{x-1}$, καθώς:
\[ g'(x) = \frac{f'(x)(x-1) - f(x) \cdot 1}{(x-1)^2} \]
Συνεπώς, έχουμε $g'(x) = e^x$. Ολοκληρώνοντας το διάστημα $[2, x]$ προκύπτει:
\[ \int_{2}^{x} g'(t) \d t = \int_{2}^{x} e^t \d t \implies [g(t)]_2^x = [e^t]_2^x \]
\[ g(x) - g(2) = e^x - e^2 \]
Από το δεδομένο $f(2) = e^2$ έχουμε $g(2) = \frac{f(2)}{2-1} = e^2$. Άρα:
\[ g(x) - e^2 = e^x - e^2 \implies g(x) = e^x \]
Επομένως, $\frac{f(x)}{x-1} = e^x$, οπότε ο τύπος της συνάρτησης είναι ο $f(x) = (x-1)e^x$.

\item Η παράγωγος της συνάρτησης είναι:
\[ f'(x) = 1 \cdot e^x + (x-1)e^x = e^x(1 + x - 1) = xe^x \]
Για $x > 1$, ισχύει $x > 0$ και $e^x > 0$, άρα $f'(x) > 0$, επομένως η συνάρτηση είναι \textbf{γνησίως αύξουσα} στο $(1, +\infty)$.
Αν διευρύνουμε το πεδίο ορισμού στο $[1, +\infty)$, η συνάρτηση παραμένει \textbf{συνεχώς} και \textbf{γνησίως αύξουσα}. Για $x=1$ έχουμε $f(1) = (1-1)e^1 = 0$.
Επειδή η συνάρτηση είναι γνησίως αύξουσα για $x \ge 1$, το ελάχιστο $y=0$ επιτυγχάνεται μόνο στο $x=1$.
Το σύνολο τιμών της συνάρτησης στο $[1, +\infty)$ είναι το $[0, +\infty)$, καθώς $\lim_{x \to +\infty} (x-1)e^x = +\infty$.

\item Θεωρούμε τη συνάρτηση $f(x) = (x-1)e^x$ στο διάστημα $[x, x+1]$ με $x > 1$. Η συνάρτηση είναι \textbf{συνεχώς} στο $[x, x+1]$ και \textbf{παραγωγήσιμη} στο $(x, x+1)$.
Σύμφωνα με το Θεώρημα Μέσης Τιμής \eng{Lagrange}, υπάρχει $\xi \in (x, x+1)$ τέτοιο ώστε:
\[ \frac{f(x+1) - f(x)}{(x+1) - x} = f'(\xi) \implies f(x+1) - f(x) = f'(\xi) \]
Αντικαθιστώντας τους τύπους:
\[ (x+1-1)e^{x+1} - (x-1)e^x = \xi e^\xi \implies x e^{x+1} - (x-1)e^x = \xi e^\xi \]
Θεωρούμε τη συνάρτηση $h(t) = te^t$ για $t > 0$. Η παράγωγός της είναι $h'(t) = e^t + te^t = e^t(1+t) > 0$, άρα η $h$ είναι \textbf{γνησίως αύξουσα}.
Επειδή $x < \xi < x+1$, ισχύει $h(\xi) > h(x)$, δηλαδή $\xi e^\xi > xe^x$.
Συνεπώς, $x e^{x+1} - (x-1)e^x > x e^x$.

\item Η εφαπτομένη $\epsilon$ της $C_f$ στο σημείο $x=2$ έχει εξίσωση:
\[ y - f(2) = f'(2)(x-2) \implies y - e^2 = 2e^2(x-2) \implies y = 2e^2 x - 4e^2 + e^2 \implies y = e^2(2x-3) \]
Η συνάρτηση $f$ είναι κυρτή στο $[1, 2]$ καθώς $f''(x) = (xe^x)' = e^x + xe^x = e^x(1+x) > 0$.
Επομένως, η εφαπτομένη $\epsilon$ βρίσκεται κάτω από την καμπύλη $C_f$ στο διάστημα $(1, 2)$. Το εμβαδόν $E$ είναι:
\[ E = \int_{1}^{2} [f(x) - \epsilon(x)] \d x = \int_{1}^{2} [(x-1)e^x - e^2(2x-3)] \d x \]
Υπολογίζουμε τα δύο ολοκληρώματα:
\begin{align*}
\int_{1}^{2} (x-1)e^x \d x &= [(x-2)e^x]_1^2 = (2-2)e^2 - (1-2)e^1 = 0 - (-e) = e \\
\int_{1}^{2} e^2(2x-3) \d x &= [e^2(x^2-3x)]_1^2 = e^2(4-6) - e^2(1-3) = -2e^2 - (-2e^2) = 0
\end{align*}
Άρα, το εμβαδόν είναι $E = e - 0 = e$ τετραγωνικές μονάδες.
\end{erwthma}
\end{thema}
%=========== ΤΕΛΟΣ ΛΥΣΗΣ - Χο ΘΕΜΑ Γ ===========